% !TeX program = xelatex
%% =====================================================================
%%  แบบฝึกหัดท้ายบทที่ 1 — ตัวบ่งปริมาณ (Quantifiers)
%%  รายวิชา 4091606 คณิตศาสตร์สำหรับคอมพิวเตอร์
%%  มหาวิทยาลัยราชภัฏอุตรดิตถ์
%% =====================================================================
\documentclass[a4paper,14pt]{article}
\usepackage{fontspec}
\usepackage{amsmath,amssymb,amsthm}
\usepackage[margin=2.5cm]{geometry}
\usepackage{xcolor}
\usepackage{enumitem}
\usepackage{array}
\usepackage{tabularx}
\usepackage{tcolorbox}
\usepackage{multicol}
\usepackage{hyperref}

\XeTeXlinebreaklocale "th_TH"
\XeTeXlinebreakskip = 0pt plus 1pt
\setmainfont[Script=Thai,Scale=1.0]{TH Sarabun New}

% Colors
\definecolor{navy}{HTML}{1F3A93}
\definecolor{Tgreen}{HTML}{2E8B57}
\definecolor{Fred}{HTML}{C0392B}
\definecolor{gold}{HTML}{B7950B}
\definecolor{lightnavy}{HTML}{E7ECF8}
\definecolor{mint}{HTML}{E6F4EA}
\definecolor{cream}{HTML}{FFF7DD}
\definecolor{lyellow}{HTML}{FFF6D5}
\definecolor{pink}{HTML}{FBE7E4}

\setlength{\parindent}{0pt}
\pagestyle{empty}

% Custom commands
\newcommand{\TT}{\textcolor{Tgreen}{\textbf{T}}}
\newcommand{\FF}{\textcolor{Fred}{\textbf{F}}}
\newcommand{\A}{\textcolor{navy}{$\forall$}}
\newcommand{\E}{\textcolor{gold}{$\exists$}}
\newcommand{\notset}{\(\forall\,x\)}

\begin{document}

%% ===== Header =====
\begin{tcolorbox}[colback=navy,colframe=navy,coltext=white,
  arc=4mm,boxrule=0pt,left=10pt,right=10pt,top=8pt,bottom=8pt]
\begin{center}
{\Large\bfseries บทที่ 1 — ตัวบ่งปริมาณ (Quantifiers)}\\[2pt]
{\large แบบฝึกหัดท้ายบทเรียน · ส่งสัปดาห์ถัดไป}\\[2pt]
{\normalsize รายวิชา 4091606 คณิตศาสตร์สำหรับคอมพิวเตอร์ · มหาวิทยาลัยราชภัฏอุตรดิตถ์}
\end{center}
\end{tcolorbox}

\vspace{0.4em}

%% ===== Student info =====
\noindent
\begin{tabularx}{\textwidth}{@{}l X l X@{}}
ชื่อ-สกุล: & \dotfill & รหัส: & \dotfill \\
กลุ่ม: & \dotfill & วันที่ส่ง: & \dotfill \\
\end{tabularx}

\vspace{0.3em}

%% ===== Instructions =====
\begin{tcolorbox}[colback=cream,colframe=gold,arc=3mm,boxrule=0.5pt,
  left=10pt,right=10pt,top=6pt,bottom=6pt]
\textbf{คำชี้แจง:}
\begin{itemize}[leftmargin=*,topsep=2pt,itemsep=1pt]
  \item ทำลงในกระดาษ A4 แสดงวิธีทำให้ชัดเจน
  \item แบบฝึกหัดแบ่งเป็น \textbf{3 ส่วน} (Part A, B, C) รวม 10 ข้อ
  \item คะแนนรวม \textbf{20 คะแนน} (Part A: 6 + Part B: 9 + Part C: 7 + โบนัส: 2)
\end{itemize}
\end{tcolorbox}

\vspace{0.5em}

%% ===== Part A =====
\begin{tcolorbox}[colback=lightnavy,colframe=navy,arc=3mm,boxrule=0.5pt,
  left=10pt,right=10pt,top=6pt,bottom=6pt]
\textbf{\textcolor{navy}{Part A · พื้นฐาน}} \hfill \textbf{2 คะแนน/ข้อ} · \textbf{6 คะแนน}
\end{tcolorbox}

\begin{enumerate}[label=\textbf{A\arabic*.},leftmargin=2em,itemsep=10pt]
  \item เขียนประโยคต่อไปนี้ให้อยู่ในรูปสัญลักษณ์ตรรกศาสตร์\\
        ``\textbf{ทุกจังหวัด}ในประเทศไทยมีผู้ว่าราชการจังหวัด''
        \vspace{1.5em}

  \item เขียนประโยคต่อไปนี้ให้อยู่ในรูปสัญลักษณ์ตรรกศาสตร์\\
        ``\textbf{มีนักเรียนบางคน}ในชั้นเรียนที่\textbf{ไม่ชอบ}วิชาคณิตศาสตร์''
        \vspace{1.5em}

  \item เขียนประโยคต่อไปนี้ให้อยู่ในรูปสัญลักษณ์ตรรกศาสตร์\\
        ``\textbf{ไม่มี}นักเรียนคนใดในห้อง\textbf{เกิดเดือนกุมภาพันธ์}''

        \emph{คำแนะนำ: ใช้ \(\forall\) ได้ เพราะ ``ไม่มี'' = ``ทุกตัวไม่เป็น''}
        \vspace{1.5em}
\end{enumerate}

\vspace{0.5em}

%% ===== Part B =====
\begin{tcolorbox}[colback=mint,colframe=Tgreen,arc=3mm,boxrule=0.5pt,
  left=10pt,right=10pt,top=6pt,bottom=6pt]
\textbf{\textcolor{Tgreen}{Part B · ตรวจสอบค่าความจริง}} \hfill \textbf{3 คะแนน/ข้อ} · \textbf{9 คะแนน}
\end{tcolorbox}

\begin{enumerate}[label=\textbf{B\arabic*.},leftmargin=2em,itemsep=10pt]
  \item จงหาค่าความจริงของ
        \(\A x \in \{1, 2, 3\}\; [x^2 \le 9]\)

        \emph{แสดงการตรวจสอบทีละค่า}
        \vspace{2em}

  \item จงหาค่าความจริงของ
        \(\E x \in \{0, 5, 10\}\; [x \text{ เป็นจำนวนคี่}]\)

        \emph{ระบุตัวอย่างยืนยัน หรือ ตัวอย่างค้าน}
        \vspace{2em}

  \item ตรวจสอบว่า
        \(\A x \in \mathbb{N}\; [x + 1 > x]\)
        เป็นจริงหรือเท็จ พร้อมให้เหตุผล

        \vspace{2em}
\end{enumerate}

\vspace{0.5em}

%% ===== Part C =====
\begin{tcolorbox}[colback=lyellow,colframe=gold,arc=3mm,boxrule=0.5pt,
  left=10pt,right=10pt,top=6pt,bottom=6pt]
\textbf{\textcolor{gold}{Part C · ประยุกต์}} \hfill \textbf{2.33 คะแนน/ข้อ} · \textbf{7 คะแนน} · \textbf{โบนัส +2}
\end{tcolorbox}

\begin{enumerate}[label=\textbf{C\arabic*.},leftmargin=2em,itemsep=10pt]
  \item เขียนประโยคต่อไปนี้ให้อยู่ในรูปสัญลักษณ์ตรรกศาสตร์ โดยใช้ทั้งตัวบ่งปริมาณและตัวเชื่อม:
        \begin{itemize}[leftmargin=*,itemsep=2pt,topsep=2pt]
          \item ให้ \(A\) = นักศึกษาเรียนวิชา CS101
          \item ให้ \(S\) = นักศึกษาส่งงาน
          \item ให้ \(G\) = นักศึกษาได้เกรด A
        \end{itemize}
        ประโยค: ``\textbf{นักศึกษาทุกคนที่เรียน CS101 และส่งงาน จะได้เกรด A}''
        \vspace{2em}

  \item เขียนประโยคต่อไปนี้ให้อยู่ในรูปสัญลักษณ์ตรรกศาสตร์ โดยใช้ตัวบ่งปริมาณซ้อน:
        \begin{itemize}[leftmargin=*,itemsep=2pt,topsep=2pt]
          \item ให้ \(M(x,y)\) = นักศึกษา \(x\) ชอบวิชา \(y\)
        \end{itemize}
        ประโยค: ``\textbf{มีนักศึกษาบางคน}ที่ชอบทั้ง\textbf{คณิตศาสตร์และเขียนโปรแกรม}''
        \vspace{2em}

  \item เขียนประโยคต่อไปนี้ให้อยู่ในรูปสัญลักษณ์ตรรกศาสตร์:
        \begin{itemize}[leftmargin=*,itemsep=2pt,topsep=2pt]
          \item ให้ \(C(x)\) = คอมพิวเตอร์เครื่องที่ \(x\) ใช้งานได้
        \end{itemize}
        ประโยค: ``\textbf{ห้องสมุดทุกเครื่อง}คอมพิวเตอร์\textbf{ใช้งานได้}''

        \textbf{คำถามเพิ่ม:} ประโยคนี้ควรใช้ \(\A\) หรือ \(\E\) จึงเหมาะสมที่สุด เพราะเหตุใด
        \vspace{1.5em}

  \item[$\bigstar$] \textbf{โจทย์ท้าทาย (โบนัส +2 คะแนน)}: เขียนประโยคต่อไปนี้ให้อยู่ในรูปสัญลักษณ์ตรรกศาสตร์ โดยใช้ตัวบ่งปริมาณซ้อน:
        \begin{itemize}[leftmargin=*,itemsep=2pt,topsep=2pt]
          \item ให้ \(F(x,y)\) = นักเรียน \(x\) เป็นเพื่อนกับ \(y\)
        \end{itemize}
        ประโยค: ``\textbf{นักเรียนทุกคน}ในชั้นเรียนนี้ \textbf{มีเพื่อนอย่างน้อย 1 คน}''
        \vspace{2em}
\end{enumerate}

\vspace{1em}

%% ===== Score table =====
\begin{center}
\begin{tcolorbox}[colback=lightnavy,colframe=navy,arc=3mm,boxrule=0.5pt,
  width=0.7\textwidth,left=10pt,right=10pt,top=6pt,bottom=6pt]
\begin{center}
\textbf{\large ตารางคะแนน}\\[4pt]
\begin{tabular}{|l|c|c|c|}
\hline
\textbf{ส่วน} & \textbf{ข้อ} & \textbf{คะแนน/ข้อ} & \textbf{รวม} \\ \hline
Part A · พื้นฐาน & 3 & 2 & 6 \\ \hline
Part B · ตรวจสอบ & 3 & 3 & 9 \\ \hline
Part C · ประยุกต์ & 3 & 2.33 & 7 \\ \hline
\textbf{รวม (ข้อบังคับ)} & \textbf{9} & --- & \textbf{20} \\ \hline
โบนัส (ท้าทาย) & 1 & 2 & 2 \\ \hline
\textbf{รวมทั้งหมด (เต็ม)} & \textbf{10} & --- & \textbf{22} \\ \hline
\end{tabular}
\end{center}
\end{tcolorbox}
\end{center}

\vspace{1em}

%% ===== Reference =====
\begin{tcolorbox}[colback=white,colframe=line,arc=2mm,boxrule=0.3pt,
  left=8pt,right=8pt,top=4pt,bottom=4pt]
\begin{center}
\textbf{อ้างอิงสัญลักษณ์}\\
\(\A x\,P(x)\) = ``สำหรับทุก \(x\), \(P(x)\) เป็นจริง'' \quad
\(\E x\,P(x)\) = ``มี \(x\) บางตัวที่ทำให้ \(P(x)\) เป็นจริง''
\end{center}
\end{tcolorbox}

\vfill

\begin{center}
\small\itshape\color{gray}
``คณิตศาสตร์สำหรับคอมพิวเตอร์'' · บทที่ 1 ตรรกศาสตร์ · สัปดาห์ที่ 3
\end{center}

\end{document}
